\section{Geometry Diagnostics}
\label{sec:geometry}

Phase~8 began with probes on the previous Phase~5 checkpoint before any
new training objective was introduced.  The purpose was to measure the
geometry that explicit delta prediction would operate in.

\begin{table}[t]
\centering
\caption{Phase~8 geometry diagnostics on a stabilized Phase~5
checkpoint.  These measurements motivate the move from direct final-layer
delta prediction to intermediate-loop supervision.}
\label{tab:geometry}
\smallskip
\begin{tabular}{lll}
\toprule
Metric & Value & Interpretation \\
\midrule
Copy cosine & $0.9999$ & Consecutive target states nearly identical \\
Final delta/state ratio & $0.010$ & Final-layer delta is about 1\% signal \\
Loop 1 delta/state ratio & $0.123$ & Early loop has about $12\times$ more signal \\
Straightness score & $0.132$ & Trajectories are mostly curved or random \\
Consecutive delta cosine & $-0.059$ & Adjacent delta directions are uncorrelated \\
Effective rank, states & $5.1/128$ & State space is highly anisotropic \\
Effective rank, deltas & $4.9/128$ & Deltas occupy the same low-rank subspace \\
Top singular direction & $64\%$ state, $85\%$ target & One direction dominates \\
Linear delta probe cosine & $0.30$ & Weak ceiling for simple delta prediction \\
\bottomrule
\end{tabular}
\end{table}

\subsection{Signal Decays Across Loops}

The most important diagnostic is the layer-wise
transition-to-state norm ratio:
\begin{center}
\begin{tabular}{lcccc}
\toprule
Loop & 1 & 2 & 3 & 6 \\
\midrule
$\|\Delta z\|/\|z\|$ & $0.123$ & $0.055$ & $0.028$ & $0.010$ \\
\bottomrule
\end{tabular}
\end{center}

The final representation is not merely a deeper version of the first
loop.  It is a representation in which most of the transition signal has
been compressed away.  Edit-type probes show the same pattern: the
discriminative information peaks around early or middle loops and drops
at the final readout.

\subsection{Why Final-Layer Deltas Are a Bad Target}

Direct delta prediction assumes that $z_{t+1}-z_t$ is a stable,
learnable object.  The probes show the opposite in this regime.  The
delta is tiny relative to the state, adjacent deltas are not aligned,
and both states and deltas lie in a low-rank anisotropic space.  In that
setting, a predictor can spend most of its capacity matching magnitude
and orientation artifacts rather than learning a reusable transition
operator.

The paper's central claim follows from this geometry: the useful
transition signal lives earlier than the final loop, and it must be
supervised before the shared block smooths it away.
